% !TeX root = ../report.tex
\section{Discussion}
\label{sec:discussion}

Taken together, the experiments show how the model's parts interact. Dust-heating feedback creates a threshold for storm growth, and enough available surface dust removes that threshold. Each global event depletes the reservoir that supplies its lifting, and slow replenishment sets when a later season can produce another event. In the baseline comparison, variable forcing turns the resulting two-year schedule into irregular spacing, much as a low-order circulation model produced year-to-year variability in global storms only once random forcing was added~\cite{pankine2002}.

These results describe the model, not Mars. The seasonal comparison covers one baseline parameter set, so it cannot show whether randomness is required for irregular spacing elsewhere, and no model component was removed to test whether it was necessary. Without geography, the model identifies ``global'' events by dust level rather than spatial extent. Most parameters are illustrative rather than fitted, so the 29.4\% baseline rate is not a test against the observed rate of about 25\%. The simulated records are finite, and year-to-year variability enters only as a uniform multiplier on the forcing.

The model resembles seasonally forced epidemic models, in which epidemics deplete susceptible people, births replenish them, and changes in the birth rate shift measles between annual, biennial, and irregular epidemics~\cite{earn2000}; the resemblance is structural only.

The MATLAB code and LaTeX source for this report are available at
\url{https://github.com/lkaising/mars-dust-project}.
